\documentclass[12pt,letterpaper]{article}

\usepackage[utf8]{inputenc}
\usepackage[T1]{fontenc}
\usepackage[spanish,es-tabla]{babel}
\usepackage[margin=1in]{geometry}
\usepackage{setspace}
\usepackage{fancyhdr}
\usepackage{enumitem}
\usepackage{hyperref}
\usepackage{titlesec}

\hypersetup{colorlinks=true, linkcolor=black, urlcolor=black, citecolor=black,
            breaklinks=true}
\usepackage{url}
\Urlmuskip=0mu plus 1mu\relax
\makeatletter
\g@addto@macro{\UrlBreaks}{\UrlOrds}
\makeatother
\urlstyle{same}
\makeatletter
\makeatother

\onehalfspacing

\pagestyle{fancy}
\fancyhf{}
\fancyhead[L]{\small Pulso - Predice, nutre, entrena y protege tu corazón}
\fancyhead[R]{\small\thepage}
\renewcommand{\headrulewidth}{0pt}

\titleformat{\section}{\normalfont\large\bfseries}{\thesection.}{0.6em}{\MakeUppercase}
\titleformat{\subsection}{\normalfont\normalsize\bfseries}{\thesubsection}{0.6em}{}

\setlist[itemize]{itemsep=4pt, topsep=6pt}
\setlist[enumerate]{itemsep=6pt, topsep=6pt}

\begin{document}

%======================= PORTADA =======================
\begin{titlepage}
\centering
\vspace*{1cm}

{\LARGE\bfseries Pulso - Predice, nutre, entrena y protege tu corazón\par}

\vspace{1.5cm}

{\large Plataforma PWA con IA para la Prevención Cardiovascular Global:\par}
\vspace{0.4cm}
{\large Predicción de Riesgo, Alimentación, Ejercicio y Seguimiento\par}
\vspace{0.4cm}
{\large Personalizado\par}

\vspace{2.5cm}

{\large Natalia Moncaleano Montero\par}
\vspace{0.4cm}
{\large Agustín Peralta Guarín\par}
\vspace{0.4cm}
{\large Diego Fernando Fandiño Olarte\par}

\vspace{2.5cm}

{\large Universidad de San Buenaventura\par}
\vspace{0.4cm}
{\large Facultad de Ingeniería (Bogotá)\par}
\vspace{0.4cm}
{\large Ingeniería de Sistemas\par}
\vspace{0.4cm}
{\large Bogotá D.C., Colombia\par}
\vspace{0.4cm}
{\large 2026\par}

\vfill
\end{titlepage}

\setcounter{page}{1}

%======================= DEDICATORIA =======================
\section*{DEDICATORIA}
\addcontentsline{toc}{section}{Dedicatoria}

Dedicamos este trabajo a nuestras familias, por su amor incondicional, comprensión
y apoyo constante durante todo el proceso de formación profesional. A nuestros
padres, quienes con su esfuerzo, valores y ejemplo nos impulsaron a superar cada
reto y alcanzar nuestras metas.

También dedicamos este logro a nuestros docentes, por compartir su conocimiento y
guiarnos con paciencia y compromiso, y a todas las personas que creyeron en
nosotros y nos motivaron a culminar con éxito este proyecto. Este trabajo es
reflejo del esfuerzo conjunto, la perseverancia y la pasión que compartimos como
grupo.

\newpage

%======================= INDICE =======================
\renewcommand{\contentsname}{ÍNDICE}
\tableofcontents

\newpage

%======================= RESUMEN =======================
\section*{RESUMEN}
\addcontentsline{toc}{section}{Resumen}

Las enfermedades cardiovasculares (ECV) constituyen la principal causa de muerte
en el mundo. De acuerdo con la Organización Mundial de la Salud (2026a), cada año
fallecen cerca de 20,5 millones de personas por esta causa, lo que equivale a una
de cada tres muertes registradas globalmente. La carga es desigual: los países de
ingresos medios y bajos concentran la mayor mortalidad evitable, en buena parte
porque sus poblaciones no acceden a herramientas de tamizaje y orientación
preventiva.

En ese contexto, este proyecto desarrolla \textbf{Pulso}, una Aplicación Web
Progresiva (PWA) basada en inteligencia artificial que articula cinco funciones
complementarias: estimación individual del riesgo cardiovascular a partir de
indicadores de estilo de vida; recomendaciones alimenticias personalizadas con
ingredientes, cantidades y preparaciones; planes de vida saludable que integran
manejo del estrés, calidad del sueño e hidratación; rutinas de ejercicio
progresivas según nivel de riesgo y condición física; y un sistema de seguimiento
con dashboard y reportes periódicos.

La arquitectura PWA permite que Pulso sea instalable en cualquier dispositivo,
funcione offline mediante Service Workers y envíe notificaciones push para
alertas preventivas. La estimación de riesgo se sustenta en un modelo supervisado
entrenado sobre la Encuesta Nacional de Examen de Salud y Nutrición de los
Estados Unidos, ciclo agosto 2021 a agosto 2023 (Centers for Disease Control and
Prevention [CDC], 2024), con una muestra analítica de 5.043 adultos y una
prevalencia de eventos cardiovasculares del 11,8\,\%. El modelo alcanza un área
bajo la curva ROC de 0,812 al incorporar edad y sexo, frente a 0,605 utilizando
únicamente los indicadores de estilo de vida que la aplicación captura hoy
(frecuencia cardíaca en reposo, peso, horas de sueño y nivel de estrés). Los
conjuntos UCI Heart Disease (Janosi et al., 1988), Framingham Heart Study (Dawber
et al., 1951) y el Global Health Observatory de la OMS (2026b) operan como
fundamento bibliográfico y epidemiológico, según se precisa en el apartado de
limitaciones. La interpretación del resultado en lenguaje natural se delega a un
modelo de lenguaje.

\vspace{0.4cm}
\noindent\textbf{Palabras clave:} enfermedades cardiovasculares, inteligencia
artificial, PWA, estimación de riesgo, salud digital, recomendaciones
alimenticias, vida saludable, sistemas de apoyo a la decisión.

\newpage

%======================= ABSTRACT =======================
\section*{ABSTRACT}
\addcontentsline{toc}{section}{Abstract}

\selectlanguage{english}

Cardiovascular diseases (CVDs) are the leading cause of death worldwide.
According to the World Health Organization (2026a), approximately 20.5 million
people die from CVDs every year, accounting for nearly one in every three
recorded deaths. The burden is unequally distributed: low- and middle-income
countries concentrate most of the preventable mortality, largely because their
populations lack access to digital screening tools and preventive guidance.

This work presents \textbf{Pulso}, an AI-driven Progressive Web Application (PWA)
that combines five complementary functions: individual cardiovascular risk
estimation from lifestyle indicators; personalized dietary recommendations
including ingredients, quantities and cooking methods; comprehensive
healthy-lifestyle plans covering stress management, sleep quality and hydration;
exercise routines that adapt progressively to risk level and fitness; and a
continuous tracking layer with dashboard and configurable periodic reports.

The PWA architecture allows Pulso to be installable on any device, to work
offline through Service Workers and to deliver push notifications for preventive
alerts. Risk estimation is grounded in a supervised model trained on the United
States National Health and Nutrition Examination Survey, August 2021 to August
2023 cycle (Centers for Disease Control and Prevention [CDC], 2024), with an
analytic sample of 5,043 adults and an 11.8\,\% prevalence of cardiovascular
events. The model reaches an area under the ROC curve of 0.812 when age and sex
are included, compared to 0.605 using only the lifestyle indicators the
application currently captures (resting heart rate, weight, sleep hours and
stress level). The UCI Heart Disease (Janosi et al., 1988), Framingham Heart
Study (Dawber et al., 1951) and WHO Global Health Observatory (2026b) datasets
serve as bibliographic and epidemiological grounding, as detailed in the
limitations section. Natural-language interpretation of the resulting score is
delegated to a language model.

\vspace{0.4cm}
\noindent\textbf{Keywords:} cardiovascular diseases, artificial intelligence,
PWA, risk estimation, digital health, dietary recommendations, healthy lifestyle,
decision support systems.

\selectlanguage{spanish}

\newpage

%======================= 1. INTRODUCCION =======================
\section{Introducción}

Las enfermedades cardiovasculares (ECV) son hoy uno de los problemas de salud
pública más serios a nivel global. La Organización Mundial de la Salud (2026a)
estima en 20,5 millones las muertes anuales por esta causa, cifra que equivale al
33\,\% de toda la mortalidad mundial y que viene en aumento sostenido año tras
año. En Colombia, el Ministerio de Salud y Protección Social (2022) reporta que
las enfermedades del sistema circulatorio son la primera causa de mortalidad
general, con una tasa ajustada superior a 130 muertes por cada 100.000
habitantes.

El panorama se complica por hábitos cada vez más extendidos: sedentarismo,
alimentación con exceso de sodio y grasas saturadas, tabaquismo, obesidad y altos
niveles de estrés. Estos factores no son exclusivos de países desarrollados ni de
economías emergentes: aparecen, con matices, en ambas. A esto se suma que muchas
personas ignoran su nivel de riesgo cardiovascular hasta que ocurre un evento
agudo, típicamente un infarto o un accidente cerebrovascular, lo que reduce
drásticamente las posibilidades de intervención preventiva.

Desde la ingeniería de sistemas existe una oportunidad concreta. Es posible
diseñar aplicaciones web progresivas (PWA) inteligentes que recojan datos de
estilo de vida del usuario, estimen su riesgo cardiovascular y devuelvan
recomendaciones accionables en tiempo real. La arquitectura PWA fue elegida
precisamente porque ofrece la experiencia de una app nativa en cualquier
navegador, sin pasar por tiendas de aplicaciones, con soporte offline y
notificaciones push para alertas de salud.

La literatura científica ha venido respaldando este enfoque. Weng et al. (2017)
demostraron que los modelos de aprendizaje automático pueden superar a los
algoritmos clínicos clásicos en la predicción de eventos cardiovasculares;
Krittanawong et al. (2020), en su meta-análisis sobre aprendizaje automático en
medicina cardiovascular, concluyen que las redes neuronales profundas y los
algoritmos de ensemble learning alcanzan exactitudes superiores al 85\,\% en la
predicción de eventos cardíacos a partir de datos clínicos estructurados. Es
importante señalar que esos desempeños se obtienen sobre variables clínicas de
laboratorio, condición que delimita el alcance del presente prototipo y que se
cuantifica en la Sección 4.3. Pulso integra esa evidencia en una PWA
que no solo estima riesgo: también recomienda alimentación con ingredientes
específicos, planes de vida saludable y rutinas de ejercicio, apoyándose en datos
epidemiológicos mundiales.

El presente trabajo revisa las limitaciones de las soluciones actuales y propone
un enfoque viable para construir una PWA inteligente de alcance global, con
recomendaciones contextualizadas por región, cultura alimentaria y disponibilidad
local de ingredientes.

\newpage

%======================= 2. PLANTEAMIENTO =======================
\section{Planteamiento del Problema}

Las enfermedades cardiovasculares son hoy el principal desafío de salud pública
tanto a escala global como en Colombia. Según la OMS (2026a), causan cerca de
20,5 millones de muertes al año; las manifestaciones más frecuentes son
hipertensión arterial, cardiopatía isquémica y accidente cerebrovascular, todas
con tendencia ascendente en la última década.

Los factores de riesgo más prevalentes, sedentarismo, dietas con alto sodio y
grasas saturadas, tabaquismo, obesidad y estrés crónico, atraviesan economías de
todo nivel de ingreso. El problema operativo es que una gran parte de la
población no conoce su riesgo cardiovascular hasta que aparece un evento clínico,
porque no existen herramientas digitales gratuitas, accesibles y de alcance
masivo que permitan una evaluación preventiva personalizada (Weng et al., 2017).

Desde la ingeniería de sistemas se identifica una brecha tecnológica clara: las
plataformas existentes suelen estar diseñadas para entornos hospitalarios, no
funcionan offline, no se instalan como aplicación nativa y no ofrecen
funcionalidades complementarias como recomendaciones alimenticias con
ingredientes y cantidades, planes integrales de vida saludable, rutinas de
ejercicio adaptadas al riesgo, o reportes periódicos accesibles desde cualquier
dispositivo y sin conexión.

\subsection*{Pregunta de Investigación}

A partir de lo anterior, la pregunta que orienta este trabajo es:

\begin{quote}
¿De qué manera una Aplicación Web Progresiva impulsada por inteligencia
artificial puede contribuir a la predicción, prevención y reducción del riesgo
cardiovascular a nivel mundial mediante recomendaciones personalizadas de
alimentación, ejercicio y vida saludable?
\end{quote}

\subsection*{Antecedentes}

La aplicación de inteligencia artificial al área cardiovascular cuenta con
evidencia sólida. Weng et al. (2017) desarrollaron modelos de machine learning
que predicen eventos cardiovasculares con mayor precisión que SCORE o el
Framingham Risk Score. Krittanawong et al. (2020) confirman que las redes
neuronales profundas y los métodos de ensemble superan el 85\,\% de exactitud
sobre datos clínicos estructurados.

En prescripción de ejercicio, Liao et al. (2020) reportan que los sistemas de
recomendación basados en aprendizaje automático mejoran significativamente la
adherencia del paciente frente a planes genéricos. En nutrición computacional,
Zeevi et al. (2015) demostraron que modelos de IA pueden estimar respuestas
metabólicas individuales y proponer dietas ajustadas con resultados clínicamente
significativos.

En cuanto a la arquitectura, Biduski et al. (2020) muestran que las PWA presentan
ventajas medibles sobre las apps nativas en contextos de salud pública: mejor
accesibilidad sin pasar por una store, funcionamiento offline para zonas con
conectividad intermitente, actualizaciones transparentes y un consumo de
almacenamiento marcadamente menor.

En el caso colombiano, Pinzón et al. (2020) y los informes del Observatorio
Nacional de Salud (2021) han documentado barreras estructurales como la falta de
tamizaje digital para población general, desigualdades regionales y baja adopción
tecnológica en atención primaria. El conjunto de estos antecedentes refuerza la
pertinencia de combinar IA con interfaces web progresivas en prevención
cardiovascular.

\newpage

%======================= 3. JUSTIFICACION =======================
\section{Justificación}

\subsection{Relevancia Social}

Las enfermedades cardiovasculares son también, y sobre todo, un indicador de
desigualdad. Millones de personas en todo el mundo no cuentan con servicios de
salud preventivos, y eso eleva la mortalidad por eventos que podrían haberse
evitado. Pulso, como PWA gratuita y de alcance global, ayuda a reducir esa
brecha: pone al alcance de cualquiera, con o sin conexión estable, una estimación
orientativa de su riesgo y una guía nutricional con ingredientes adaptados a su
contexto.

La elección de la arquitectura PWA es estratégica para países de ingresos medios
y bajos, donde la conectividad es intermitente: los usuarios pueden consultar
recomendaciones, rutinas y niveles de riesgo aun en zonas rurales o con señal
limitada, gracias al funcionamiento offline.

\subsection{Relevancia Teórica}

El proyecto se apoya en una literatura ya establecida sobre IA aplicada a
estratificación del riesgo y personalización de intervenciones preventivas (Weng
et al., 2017; Krittanawong et al., 2020). Desde la ingeniería de sistemas, este
trabajo aporta a la base teórica del uso de PWA inteligentes integradas con
sistemas de información en salud, contribuyendo a campos como la salud digital,
la nutrición computacional y la medicina preventiva de precisión.

Adicionalmente, el trabajo documenta una discusión metodológica poco frecuente en
la literatura de prototipos de salud digital: cuánta capacidad predictiva puede
alcanzarse con las variables que un usuario no clínico aporta sin supervisión
médica. La cuantificación presentada en la Sección 4.3 muestra que los
indicadores de estilo de vida por sí solos discriminan de forma limitada, y que
la ganancia decisiva proviene de dos variables demográficas de bajo costo de
recolección. Ese resultado es transferible a cualquier herramienta preventiva de
acceso abierto y constituye un aporte metodológico del proyecto.

\subsection{Relevancia Práctica}

Pulso permitirá a usuarios de cualquier país ingresar sus indicadores de estilo
de vida para obtener:

\begin{itemize}
  \item Estimación orientativa de su nivel de riesgo cardiovascular, con
    umbrales sustentados en literatura clínica y datos epidemiológicos
    mundiales.
  \item Recomendaciones alimenticias detalladas con ingredientes, cantidades,
    métodos de preparación y alternativas regionales según cultura alimentaria.
  \item Planes de vida saludable que integran manejo del estrés, calidad del
    sueño, hidratación y hábitos cotidianos.
  \item Rutinas de ejercicio adaptadas a su condición física, nivel de riesgo y
    objetivos personales.
  \item Alertas preventivas automáticas mediante notificaciones push.
  \item Reportes periódicos configurables (mensuales o quincenales), accesibles
    incluso sin conexión, que muestren la evolución de sus indicadores.
\end{itemize}

\newpage

%======================= 4. ALCANCE Y LIMITACIONES =======================
\section{Alcance y Limitaciones}

\subsection{Alcance}

El proyecto contempla el desarrollo de Pulso como una PWA de alcance global,
basada en inteligencia artificial, orientada a la estimación del riesgo
cardiovascular, la generación de recomendaciones alimenticias con ingredientes
específicos, planes de vida saludable, rutinas de ejercicio personalizadas y el
seguimiento continuo del estado de salud del usuario.

La plataforma estará disponible para cualquier persona con acceso a internet y
navegador compatible, será instalable en móvil, tablet o desktop, funcionará
offline mediante Service Workers y enviará notificaciones push para alertas
preventivas. El sistema integra seis módulos funcionales:

\begin{enumerate}
  \item Ingreso de indicadores de estilo de vida asociado a un identificador
    anónimo generado en el propio dispositivo, sin registro de usuario ni
    recolección de datos personales identificables.

  \item Estimación del riesgo cardiovascular a partir de frecuencia cardíaca en
    reposo, peso, horas de sueño y nivel de estrés percibido. El prototipo
    implementa actualmente un modelo heurístico ponderado, cuyos umbrales se
    sustentan en la literatura clínica, y el trabajo incorpora el entrenamiento y
    la evaluación de un modelo supervisado sobre datos poblacionales reales
    (Sección 4.3) que constituye la base del estimador definitivo. La
    explicabilidad es directa por construcción: tanto el modelo heurístico como
    la regresión logística informan la contribución de cada indicador sobre el
    puntaje final, sin requerir métodos de atribución post hoc.

  \item Recomendaciones alimenticias personalizadas con ingredientes, cantidades,
    métodos de preparación, alternativas regionales y planes semanales con lista
    de compras automatizada.

  \item Recomendaciones de vida saludable que abarcan manejo del estrés, calidad
    del sueño, hidratación y hábitos cotidianos.

  \item Rutinas de ejercicio adaptadas al nivel de riesgo con progresividad
    automática y alertas de seguridad.

  \item Sistema de seguimiento con dashboard interactivo, alertas preventivas y
    reportes periódicos exportables en PDF.
\end{enumerate}

\subsection{Limitaciones}

El proyecto presenta varias limitaciones que conviene declarar de entrada. El
desarrollo se limita a un prototipo funcional en fase experimental, sin
implementación directa en instituciones médicas, y la validación se centra en
métricas computacionales de desempeño y en pruebas de usabilidad.

Las recomendaciones generadas son de carácter orientativo y no sustituyen el
criterio del profesional de salud. La disponibilidad de ingredientes puede variar
según la región del usuario; las funcionalidades offline dependen del soporte del
navegador para Service Workers; y la exactitud de la estimación está condicionada
por la calidad y veracidad de los datos que ingrese el usuario, dado que el
sistema no cuenta con verificación clínica directa.

Sobre los conjuntos de datos conviene precisar el alcance real de su uso. La
fuente de entrenamiento es NHANES 2021-2023 (CDC, 2024), seleccionada porque es
el único conjunto de acceso público que mide las mismas variables que la
aplicación captura. Los tres conjuntos citados inicialmente en la propuesta
fueron verificados y descartados como corpus de entrenamiento por razones
distintas, que se documentan a continuación.

El UCI Heart Disease (Janosi et al., 1988) es de acceso abierto y fue verificado:
su subconjunto usable, \texttt{processed.cleveland.data}, contiene 303 registros
y 14 variables, de los cuales 6 presentan valores faltantes. Sus predictores
exigen exámenes de laboratorio y prueba de esfuerzo (colesterol sérico, glucemia
en ayunas, electrocardiograma en reposo y frecuencia cardíaca máxima bajo
esfuerzo) que Pulso no solicita por decisión de diseño. La intersección entre sus
variables y las que captura la aplicación es nula, de modo que se conserva como
antecedente metodológico y no como fuente de datos.

El Framingham Heart Study (Dawber et al., 1951) no es de descarga libre. Su
acceso se tramita ante el repositorio BioLINCC del National Heart, Lung, and
Blood Institute y exige aprobación de un comité de ética institucional, completa
o expedita, además de un acuerdo formal de uso de datos, plazos incompatibles con
el cronograma del presente trabajo. Por su parte, el Global Health Observatory de
la OMS (2026b) publica indicadores agregados por país y no registros
individuales, por lo que sustenta las cifras epidemiológicas de contexto pero no
permite entrenar un estimador individual.

El uso de NHANES impone cuatro restricciones que condicionan la interpretación de
los resultados de la Sección 4.3 y que se declaran de manera explícita. Primera:
el diseño de la encuesta es transversal, de modo que el modelo estima la
asociación entre el estado actual del participante y el antecedente de un evento
cardiovascular, y no la incidencia futura de dicho evento; el prototipo no
realiza, por tanto, predicción prospectiva. Segunda: la variable de estrés se
aproximó mediante el cuestionario PHQ-9, instrumento validado que mide síntomas
depresivos y no estrés percibido, dado que el ciclo no incorpora una escala de
estrés; se trata de una equivalencia aproximada. Tercera: el desenlace es
autorreportado por el participante y no verificado contra historia clínica.
Cuarta: se omitieron los ponderadores de muestreo complejo, admisible para un
clasificador demostrativo pero incompatible con la generación de estimaciones
poblacionales.

Finalmente, la validación externa del estimador en población colombiana queda
planteada como línea de trabajo posterior, dado que NHANES describe a la
población de los Estados Unidos y la transferencia de un modelo de riesgo entre
poblaciones requiere recalibración.

\subsection{Resultados Preliminares del Estimador}

Con el fin de sustentar empíricamente el modelo de estimación, se construyó una
muestra analítica a partir de siete archivos del ciclo NHANES 2021-2023 (CDC,
2024), cruzados por identificador de participante. Se retuvieron los adultos de
18 años o más con las cuatro variables de Pulso completas y desenlace
cardiovascular conocido, obteniendo \textbf{5.043 participantes} con
\textbf{597 casos positivos (11,8\,\%)}. La edad media es de 53,6 años, la
frecuencia cardíaca media de 70,8 lpm, el índice de masa corporal medio de 29,7 y
la duración media del sueño de 7,7 horas.

La correspondencia entre las variables de la aplicación y las de la encuesta es
directa: la frecuencia cardíaca en reposo se toma de \texttt{BPXOPLS1}, medida
por oscilometría tras cinco minutos de reposo en posición sentada; el peso de
\texttt{BMXWT}, medido en el centro de examen y no autorreportado; las horas de
sueño de \texttt{SLD012}; y el estrés se aproxima mediante el puntaje total del
PHQ-9. El desenlace agrupa infarto, angina, enfermedad coronaria, falla cardíaca
y accidente cerebrovascular (\texttt{MCQ160B} a \texttt{MCQ160F}).

Se evaluaron cuatro modelos supervisados mediante partición estratificada 70/30 y
validación cruzada de cinco pliegues, contrastados contra la heurística ponderada
que implementa el prototipo. La Tabla 1 resume el desempeño.

\vspace{0.3cm}
\noindent\textbf{Tabla 1.} \textit{Desempeño discriminativo de los modelos
evaluados sobre NHANES 2021-2023.}

\vspace{0.2cm}
\noindent\begin{tabular}{p{5.1cm}p{3.6cm}cc}
\hline
\textbf{Modelo} & \textbf{Variables} & \textbf{AUC} & \textbf{AUC VC} \\
\hline
Heurística ponderada actual & FC, sueño, estrés & 0,565 & -- \\
Regresión logística & Las 4 de la app & 0,605 & 0,600 \\
Gradient boosting & Las 4 de la app & 0,572 & 0,590 \\
Regresión logística & Las 4 + edad y sexo & \textbf{0,812} & 0,802 \\
Gradient boosting & Las 4 + edad y sexo & 0,787 & 0,787 \\
\hline
\end{tabular}

\vspace{0.2cm}
\noindent\footnotesize AUC: área bajo la curva ROC sobre la partición de prueba
(n = 1.513). AUC VC: media de validación cruzada de cinco pliegues. La heurística
se evaluó invirtiendo el puntaje, dado que en la aplicación un valor alto indica
mejor estado.\normalsize

\vspace{0.4cm}

El resultado principal es que los cuatro indicadores de estilo de vida que la
aplicación captura poseen, por sí solos, una capacidad discriminativa limitada:
un AUC de 0,605 se sitúa apenas por encima del azar, y la heurística vigente
alcanza 0,565. La incorporación de edad y sexo eleva el AUC a 0,812, con una
sensibilidad de 0,793 y una especificidad de 0,702 en el punto de Youden.

Los odds ratio estandarizados explican el fenómeno: la edad presenta un OR de
4,16 por desviación estándar y el sexo de 0,64, magnitudes que dominan sobre el
índice de masa corporal (1,73), el PHQ-9 (1,52) y la frecuencia cardíaca (1,15).
Este hallazgo tiene una consecuencia directa de diseño, y es que el factor
limitante del estimador no es la ausencia de variables de laboratorio, como se
presumía en la formulación inicial del proyecto, sino la ausencia de dos campos
demográficos que el usuario puede aportar sin intervención clínica. En
consecuencia, se incorpora al alcance la solicitud de edad y sexo en el registro
inicial de la aplicación.

Dos precisiones sobre los coeficientes. El peso y el índice de masa corporal son
colineales, por lo que el modelo reparte el efecto entre ambos y el peso aparece
con un OR inferior a la unidad, situación que no debe leerse como un efecto
protector. De igual forma, la duración del sueño presenta un OR de 1,12, atribuible
a confusión por edad y morbilidad previa y no a una relación causal.

\newpage

%======================= 5. OBJETIVOS =======================
\section{Objetivos}

\subsection{Objetivo General}

Desarrollar Pulso, una Aplicación Web Progresiva (PWA) basada en inteligencia
artificial para la estimación del riesgo cardiovascular, la generación de
recomendaciones alimenticias con ingredientes específicos, planes de vida
saludable, rutinas de ejercicio personalizadas y seguimiento continuo, orientada
a la prevención cardiovascular a nivel mundial.

\subsection{Objetivos Específicos}

\begin{itemize}
  \item Diseñar la arquitectura de Pulso como PWA, definiendo y justificando la
    selección tecnológica de sus componentes, la integración de Service Workers,
    persistencia local y notificaciones push, junto con el modelo de estimación
    del riesgo cardiovascular y el sustento bibliográfico de sus umbrales y
    ponderaciones.

  \item Desarrollar los módulos funcionales de Pulso: estimación de riesgo
    cardiovascular, recomendaciones alimenticias con ingredientes específicos,
    planes de vida saludable, rutinas de ejercicio con progresividad automática y
    sistema de seguimiento con dashboard y reportes en PDF.

  \item Validar el prototipo mediante la evaluación del estimador de riesgo sobre
    datos poblacionales reales, empleando métricas de discriminación y
    calibración con partición estratificada y validación cruzada, complementada
    con pruebas de usabilidad con usuarios representativos, la revisión de un
    profesional de salud y la verificación del funcionamiento en distintos
    dispositivos y condiciones de conectividad.
\end{itemize}

\newpage

%======================= 6. CONCLUSIONES =======================
\section{Conclusiones}

Las enfermedades cardiovasculares constituyen, simultáneamente, un desafío
médico, tecnológico y social. La detección tardía del riesgo, la fragmentación de
los datos clínicos, las barreras de acceso a servicios preventivos y la baja
alfabetización en salud configuran un escenario complejo que demanda soluciones
integrales y no fragmentadas.

Desde la ingeniería de sistemas, la combinación de estimación algorítmica del
riesgo, arquitecturas web progresivas y diseño centrado en el usuario ofrece una
vía concreta para acercar la prevención cardiovascular a poblaciones que hoy no
tienen acceso a herramientas digitales de tamizaje. Una PWA como Pulso permite
que la evaluación de riesgo, las recomendaciones alimenticias y los planes de
ejercicio estén disponibles offline, sean instalables sin pasar por una tienda de
aplicaciones y se actualicen de forma transparente.

La revisión de antecedentes confirma que los modelos de IA pueden superar a los
algoritmos clásicos en la estratificación del riesgo (Weng et al., 2017;
Krittanawong et al., 2020), y que las arquitecturas PWA son una vía viable para
llegar a contextos con conectividad limitada (Biduski et al., 2020). Pulso
integra estos elementos en una propuesta que va más allá de la estimación:
incorpora recomendaciones nutricionales con ingredientes específicos, planes de
vida saludable y rutinas de ejercicio adaptadas al nivel de riesgo del usuario.

El trabajo también aporta evidencia empírica sobre una pregunta de diseño que la
literatura de prototipos de salud digital suele dejar implícita: cuánto riesgo
cardiovascular puede estimarse sin datos de laboratorio. La evaluación sobre
5.043 adultos de NHANES 2021-2023 muestra que los indicadores de estilo de vida
autorreportados alcanzan un AUC de 0,605, mientras que agregar edad y sexo lo
eleva a 0,812. El factor limitante no era, entonces, la ausencia de exámenes
clínicos, sino la de dos campos demográficos elementales. Reconocer y cuantificar
esa restricción, en lugar de sortearla, permite delimitar con honestidad el
alcance del prototipo y orienta su rediseño.

Como trabajo de grado, este proyecto aporta tanto al campo de la salud digital
como a la ingeniería de sistemas aplicada al sector salud, y se alinea con los
Objetivos de Desarrollo Sostenible vinculados a salud y bienestar, reducción de
desigualdades e innovación tecnológica. Las próximas fases del trabajo abordarán
el diseño metodológico detallado, los instrumentos de recolección y la validación
con usuarios representativos y revisión de un profesional de la salud.

\newpage

%======================= REFERENCIAS =======================
\section*{REFERENCIAS}
\addcontentsline{toc}{section}{Referencias}

\begin{enumerate}[label={[\arabic*]}, leftmargin=2.2em]

\item Biduski, D., Bellei, E. A., Rodriguez, J. P. M., Zaina, L. A. M., \& De
Marchi, A. C. B. (2020). Assessing long-term user experience on a mobile health
application through an in-app embedded conversation-based questionnaire.
\textit{Computers in Human Behavior, 104}, 106169.
\url{https://doi.org/10.1016/j.chb.2019.106169}

\item Centers for Disease Control and Prevention, National Center for Health
Statistics. (2024). \textit{National Health and Nutrition Examination Survey:
August 2021 -- August 2023 data files} [Conjunto de datos]. U.S. Department of
Health and Human Services.
\url{https://wwwn.cdc.gov/nchs/nhanes/continuousnhanes/default.aspx?Cycle=2021-2023}

\item Dawber, T. R., Meadors, G. F., \& Moore, F. E. (1951). Epidemiological
approaches to heart disease: The Framingham Study. \textit{American Journal of
Public Health, 41}(3), 279--281. \url{https://doi.org/10.2105/AJPH.41.3.279}

\item Detrano, R., Janosi, A., Steinbrunn, W., Pfisterer, M., Schmid, J.,
Sandhu, S., Guppy, K., Lee, S., \& Froelicher, V. (1989). International
application of a new probability algorithm for the diagnosis of coronary artery
disease. \textit{American Journal of Cardiology, 64}(5), 304--310.
\url{https://doi.org/10.1016/0002-9149(89)90524-9}

\item Janosi, A., Steinbrunn, W., Pfisterer, M., \& Detrano, R. (1988).
\textit{Heart Disease} [Conjunto de datos]. UCI Machine Learning Repository.
\url{https://doi.org/10.24432/C52P4X}

\item Krittanawong, C., Virk, H. U. H., Bangalore, S., Wang, Z., Johnson, K. W.,
Pinotti, R., Zhang, H., Kaplin, S., Narasimhan, B., Kitai, T., Baber, U.,
Halperin, J. L., \& Tang, W. H. W. (2020). Machine learning prediction in
cardiovascular diseases: A meta-analysis. \textit{Scientific Reports, 10}(1),
16057. \url{https://doi.org/10.1038/s41598-020-72685-1}

\item Liao, P., Greenewald, K., Klasnja, P., \& Murphy, S. (2020). Personalized
HeartSteps: A reinforcement learning algorithm for optimizing physical activity.
\textit{Proceedings of the ACM on Interactive, Mobile, Wearable and Ubiquitous
Technologies, 4}(1), 1--22. \url{https://doi.org/10.1145/3381007}

\item Ministerio de Salud y Protección Social de Colombia. (2022).
\textit{Análisis de Situación de Salud (ASIS) Colombia}. Bogotá D.C.
\url{https://www.minsalud.gov.co}

\item National Heart, Lung, and Blood Institute. (2026). \textit{Framingham
Heart Study: Three generations of dedication} [Sitio web]. Recuperado el 15 de
agosto de 2026, de \url{https://www.framinghamheartstudy.org}

\item Observatorio Nacional de Salud. (2021). \textit{Informe sobre enfermedades
cardiovasculares en Colombia}. Instituto Nacional de Salud.
\url{https://www.ins.gov.co}

\item Organización Mundial de la Salud. (2026a). \textit{Cardiovascular diseases
(CVDs): Key facts}. OMS.
\url{https://www.who.int/news-room/fact-sheets/detail/cardiovascular-diseases-cvds}

\item Organización Mundial de la Salud. (2026b). \textit{Global Health
Observatory} [Repositorio de datos]. Recuperado el 15 de agosto de 2026, de
\url{https://www.who.int/data/gho}

\item Pinzón, A., Aristizábal, J. C., Cubillos-Garzón, L. A., González, C., \&
Lara-Díaz, M. F. (2020). Factores de riesgo cardiovascular en población adulta
colombiana: brechas en tamizaje y prevención. \textit{Revista Colombiana de
Cardiología, 27}(4), 263--272.

\item Weng, S. F., Reps, J., Kai, J., Garibaldi, J. M., \& Qureshi, N. (2017).
Can machine-learning improve cardiovascular risk prediction using routine
clinical data? \textit{PLoS ONE, 12}(4), e0174944.
\url{https://doi.org/10.1371/journal.pone.0174944}

\item Zeevi, D., Korem, T., Zmora, N., Israeli, D., Rothschild, D., Weinberger,
A., Ben-Yacov, O., Lador, D., Avnit-Sagi, T., Lotan-Pompan, M., Suez, J., Mahdi,
J. A., Matot, E., Malka, G., Kosower, N., Rein, M., Zilberman-Schapira, G.,
Dohnalová, L., Pevsner-Fischer, M., \ldots{} Segal, E. (2015). Personalized
nutrition by prediction of glycemic responses. \textit{Cell, 163}(5),
1079--1094. \url{https://doi.org/10.1016/j.cell.2015.11.001}

\end{enumerate}

\end{document}
